\documentclass[10pt,a4paper]{report}

\usepackage[portuguese]{babel}
\usepackage{graphicx}
\usepackage{hyperref}
\usepackage{fancyhdr}

\pagestyle{fancy}
\fancyhf{}
\fancyhead[LE,RO]{\bfseries\thepage}
\fancyhead[LO]{\bfseries\rightmark}
\fancyhead[RE]{\bfseries\leftmark}
\renewcommand{\headrulewidth}{0.5pt}
\setlength{\headheight}{15.6pt}

\begin{document}

\thispagestyle{empty}
\setcounter{page}{-1}

\begin{center}
\begin{Huge}
\textbf{Universidade da Beira Interior}
\end{Huge}
\end{center}

\begin{center}
\begin{Huge}
Departamento de Informática
\end{Huge}
\end{center}

\vspace{0.1cm}
\begin{figure}[!htb]
\centering
\includegraphics[width=191pt]{ubi-fe-di.png}
\end{figure}

\vspace{0.4cm}
\begin{center}
\begin{Large}
\textbf{\emph{AutoMercado - Sistemas Distribuídos}}\\
\textbf{Licenciatura em Engenharia Informática}
\end{Large}
\end{center}

\vspace{0.4cm}
\begin{center}
\textbf{Gabriel Maria n.º 47922}\\
\textbf{Tiago Almeida n.º 48278}
\end{center}

\vspace{0.4cm}
\begin{center}
\today
\end{center}

\clearpage
\tableofcontents
\clearpage

\chapter{Introdução}

Este relatório descreve o projeto \textit{AutoMercado}, desenvolvido no âmbito da unidade curricular de Sistemas Distribuídos. A aplicação consiste num backend para uma plataforma de compra e venda de veículos, permitindo gerir utilizadores, produtos, categorias, endereços e encomendas.

A solução foi implementada em Java com Spring Boot e disponibiliza uma API REST, bem como uma interface Web simples para demonstração do funcionamento do sistema.

\chapter{Modelo de dados}

O modelo de dados da aplicação é composto pelas entidades principais:
\begin{itemize}
    \item \texttt{UtilizadorLocal}, que representa os utilizadores do sistema;
    \item \texttt{Endereco}, associado aos utilizadores;
    \item \texttt{Categoria}, usada para organizar os produtos;
    \item \texttt{Produto}, que representa os veículos ou artigos disponíveis;
    \item \texttt{Inventario}, responsável pelo controlo de stock;
    \item \texttt{Encomenda}, que regista as compras efetuadas;
    \item \texttt{QuantidadeEncomendas}, que representa os itens de cada encomenda.
\end{itemize}

As relações principais permitem associar vários endereços a um utilizador, vários produtos a uma categoria e vários itens a uma encomenda. Desta forma, o modelo suporta o fluxo completo de compra e gestão de catálogo.

\chapter{Descrição da arquitetura e das funcionalidades da aplicação}

A aplicação segue uma arquitetura em camadas, composta por controladores REST, serviços, repositórios JPA e base de dados relacional. Esta organização permite separar a lógica de negócio, o acesso a dados e a exposição dos endpoints.

Entre as principais funcionalidades implementadas destacam-se:
\begin{itemize}
    \item registo e autenticação de utilizadores;
    \item consulta de categorias e produtos;
    \item gestão de perfil e endereços;
    \item criação de encomendas com atualização de stock;
    \item consulta do histórico de compras;
    \item gestão de produtos e categorias por administradores;
    \item consulta de estatísticas de vendas.
\end{itemize}

\chapter{Descrição das opções adotadas para gestão de sessões e segurança}

A autenticação foi implementada com Spring Security. As palavras-passe são armazenadas com \texttt{BCrypt}, garantindo maior segurança no armazenamento das credenciais.

A aplicação utiliza sessões HTTP para manter o estado autenticado do utilizador entre pedidos. A autorização é feita com base em papéis, distinguindo utilizadores do tipo cliente e administrador.

Os endpoints públicos incluem login, registo e consulta de catálogo. Já as operações de perfil, compras e gestão de endereços exigem autenticação, enquanto as operações administrativas e estatísticas exigem permissões de administrador.

\chapter{Demonstração}

A demonstração da aplicação pode ser feita através de uma sequência simples de utilização. Primeiro, o utilizador consulta o catálogo sem autenticação. Em seguida, pode criar conta, iniciar sessão, adicionar um endereço e realizar uma compra.

Após a compra, é possível consultar o histórico e a fatura gerada. No caso de um utilizador administrador, ficam também disponíveis funcionalidades de gestão de produtos e categorias, bem como estatísticas de faturação e vendas.

Esta demonstração permite verificar o funcionamento global do sistema, incluindo autenticação, controlo de acessos, persistência de dados e atualização de stock.

\chapter{Conclusão}

O projeto \textit{AutoMercado} permitiu desenvolver uma aplicação funcional de comércio eletrónico orientada ao domínio automóvel, integrando persistência de dados, autenticação, controlo de acessos e gestão de encomendas.

A solução cumpre os objetivos principais do trabalho e demonstra a aplicação prática de conceitos de Sistemas Distribuídos com recurso a tecnologias como Spring Boot, Spring Data JPA e Spring Security.

\end{document}
